\chapter{Modelagem Matemática e Métodos Numéricos}
\label{cap:modelagem}

Este capítulo descreve o modelo construído para representar um canal do reator
monolítico, as hipóteses que o sustentam, a forma como a reação heterogênea foi
incorporada e os procedimentos adotados para verificar a solução numérica. As
simulações foram conduzidas no código comercial Simcenter STAR-CCM+
\parencite{StarCCMUserGuide}.

\section{Domínio e hipóteses}

O domínio computacional representa um único canal do monólito, tratado como
representativo do feixe pelas razões discutidas no
Capítulo~\ref{cap:revisao}. A Figura~\ref{fig:dominio} resume sua geometria e
as condições de contorno impostas.

Adotou-se um paralelepípedo de \SI{50}{\milli\meter} de comprimento por
\SI{1.1}{\milli\meter} de altura e \SI{50}{\micro\meter} de espessura. Nas duas
faces normais à direção de menor dimensão impôs-se condição de simetria, e a
malha contém uma única célula nessa direção. O domínio comporta-se, assim, como
um canal plano bidimensional, embora a formulação do solver permaneça
tridimensional. O artifício evita as restrições que a formulação bidimensional
estrita impõe à especificação de condições de contorno de espécies no código
empregado, sem o custo de uma terceira direção efetivamente discretizada.

A parede superior é catalítica e a inferior é tratada como parede lisa inerte.
A escolha merece registro por ser conservadora: um canal real de monólito tem
todas as paredes recobertas, de modo que a razão entre área catalítica e volume
do presente domínio é metade da de um canal plano com ambas as paredes ativas,
e um oitavo da do canal quadrado das peças efetivamente disponíveis no
laboratório. A comparação entre as duas geometrias é retomada na
Seção~\ref{sec:geomreal}; o que importa aqui é que o regime de operação
determinado no Capítulo~\ref{cap:resultados} não depende dessa razão, por ser
estabelecido localmente na parede.

\begin{figure}[htb]
\centering
\begin{tikzpicture}[x=0.78cm, y=0.78cm,
    >={Latex[length=4pt,width=3pt]},
    cota/.style={tinta3viz, line width=0.5pt, |<->|,
                 >={Latex[length=3.5pt,width=2.8pt]}},
    aux/.style={tinta3viz, line width=0.35pt, dash pattern=on 2pt off 1.5pt}]

% ---------- camada-limite de concentração (sombreado sob a parede ativa) ----
\fill[azulviz!12]
  plot[domain=0.0001:12, samples=80] ({\x}, {2.40 - 0.55*((\x/12)^(1/3))})
  -- (12,2.40) -- (0,2.40) -- cycle;
\draw[azulviz!55, line width=0.6pt]
  plot[domain=0.0001:12, samples=80] ({\x}, {2.40 - 0.55*((\x/12)^(1/3))});
\node[azulviz!75!black, font=\footnotesize, anchor=north east] at (11.7,1.70)
  {$\delta_c(x)$};

% ---------- canal -----------------------------------------------------------
\draw[tintaviz, line width=0.8pt] (0,0) rectangle (12,2.40);

% ---------- parede catalítica (superior) ------------------------------------
\fill[pattern=north east lines, pattern color=laranjaviz!65]
  (0,2.40) rectangle (12,2.70);
\draw[laranjaviz!85!black, line width=0.9pt] (0,2.40) -- (12,2.40);
\draw[laranjaviz!85!black, line width=0.9pt] (0,2.70) -- (12,2.70);
\node[font=\footnotesize, tintaviz, anchor=south] at (6,3.18)
  {parede catalítica};
\draw[cota] (12.35,2.40) -- (12.35,2.70);
\node[font=\footnotesize, tintaviz, anchor=west] at (12.5,2.55)
  {$\delta_{wc} = 10$ µm};

% ---------- parede inerte (inferior) ----------------------------------------
\fill[tinta3viz!45] (0,-0.30) rectangle (12,0);
\draw[tintaviz, line width=0.8pt] (0,-0.30) rectangle (12,0);
\node[font=\footnotesize, tintaviz, anchor=north] at (6,-0.46)
  {parede lisa inerte};

% ---------- perfil de velocidade na entrada ---------------------------------
\draw[aux] (-2.05,0) -- (-2.05,2.40);
\draw[azulviz, line width=0.9pt]
  plot[domain=0:2.40, samples=60] ({-2.05 + 1.55*4*(\x/2.40)*(1-\x/2.40)}, {\x});
\foreach \y in {0.3,0.6,0.9,1.2,1.5,1.8,2.1}{
  \draw[->, azulviz, line width=0.5pt]
    (-2.05,{\y}) -- ({-2.05 + 1.55*4*(\y/2.40)*(1-\y/2.40)}, {\y});}
\node[azulviz, font=\footnotesize, anchor=east] at (-2.2,1.20) {$\bar{u}$};

% ---------- fluxo de espécies na parede -------------------------------------
\foreach \x in {2.4,5.0,7.6}{
  \draw[->, laranjaviz!85!black, line width=0.6pt] ({\x},1.82) -- ({\x},2.32);
  \draw[->, laranjaviz!85!black, line width=0.6pt] ({\x+0.55},2.32) -- ({\x+0.55},1.82);}
\node[laranjaviz!85!black, font=\footnotesize, anchor=west] at (8.55,2.07) {$J_i$};

% ---------- rótulos de fronteira --------------------------------------------
\node[font=\footnotesize, tintaviz, anchor=north] at (0,-0.46)  {entrada};
\node[font=\footnotesize, tintaviz, anchor=north] at (12,-0.46) {saída};

% ---------- cotas ------------------------------------------------------------
\draw[aux] (0,-0.46) -- (0,-1.45);
\draw[aux] (12,-0.46) -- (12,-1.45);
\draw[cota] (0,-1.30) -- (12,-1.30);
\node[font=\footnotesize, tintaviz, fill=white, inner sep=1.5pt]
  at (6,-1.30) {$L = 50$ mm};

\draw[aux] (12,0)    -- (13.05,0);
\draw[aux] (12,2.40) -- (13.05,2.40);
\draw[cota] (12.90,0) -- (12.90,2.40);
\node[font=\footnotesize, tintaviz, anchor=west] at (13.05,1.20)
  {$H = 1{,}1$ mm};

\end{tikzpicture}
\caption{Representação esquemática do domínio computacional, fora de escala.
O escoamento entra à esquerda com perfil de Poiseuille plenamente desenvolvido;
a reação heterogênea é imposta como fluxo de espécies $J_i$ na parede superior,
recoberta pelo \textit{washcoat}. O sombreado indica a camada-limite de
concentração que se desenvolve a partir da entrada. As duas faces normais a
$z$, afastadas de \SI{50}{\micro\meter} e discretizadas com uma única célula,
recebem condição de simetria}
\label{fig:dominio}
\end{figure}

As hipóteses adotadas são as seguintes:

\begin{enumerate}
\itemsep2pt
\item escoamento laminar, incompressível e isotérmico a \SI{65}{\celsius};
\item propriedades físicas constantes, de modo que o campo de velocidade fica
      desacoplado do campo de concentrações;
\item reação restrita à superfície da parede recoberta, sem discretização da
      camada de \textit{washcoat};
\item transporte difusivo descrito pela lei de Fick com difusividades efetivas
      iguais para todas as espécies, pelas razões expostas na
      Seção~\ref{sec:difusividades};
\item cinética irreversível de pseudo-primeira ordem no triglicerídeo, com
      metanol em excesso.
\end{enumerate}

A primeira hipótese se justifica pelo número de Reynolds do escoamento, da ordem
de \num{0.2}, e pela condução do processo em banho termostatizado. A segunda
decorre das faixas estreitas de composição percorridas pelo meio reacional ao
longo do canal, dada a conversão por passe baixa que se verifica no
Capítulo~\ref{cap:resultados}. A quinta é sustentada pela razão molar de
\num{4}:\num{1} entre metanol e óleo, que mantém a concentração de álcool
praticamente inalterada ao longo da reação, e pela baixa conversão por passe,
condição sob a qual \textcite{Ezzati2021} demonstram que a aproximação de
pseudo-primeira ordem é acurada.

\section{Equações governantes}

Com massa específica e viscosidade constantes, a continuidade e a quantidade de
movimento assumem a forma

\begin{align}
\nabla \cdot \mathbf{u} &= 0 \label{eq:continuidade}\\[4pt]
\rho \left( \frac{\partial \mathbf{u}}{\partial t} + \mathbf{u} \cdot \nabla \mathbf{u} \right)
&= -\nabla p + \mu \nabla^2 \mathbf{u} \label{eq:momento}
\end{align}

O transporte das $N = 4$ espécies consideradas, a saber TG, MeOH, FAME e GL, é
descrito por

\begin{equation}
\frac{\partial (\rho Y_i)}{\partial t} + \nabla \cdot (\rho\, \mathbf{u}\, Y_i)
= \nabla \cdot \left( \rho\, D_i\, \nabla Y_i \right),
\qquad i = 1, \ldots, N
\label{eq:especies}
\end{equation}

\noindent
sem termo fonte volumétrico, uma vez que toda a reação é imposta pela condição
de contorno de parede. O solver garante $\sum_i Y_i = 1$ por normalização, o que
tem implicações discutidas na Seção~\ref{sec:difusividades}.

\section{Cinética e condição de contorno catalítica}
\label{sec:bc}

A taxa de reação por área de parede foi escrita como de primeira ordem aparente
no triglicerídeo, corrigida pelo fator de efetividade que contabiliza a
resistência interna ao \textit{washcoat}:

\begin{equation}
r_w = \eta\, k_{cat}\, w\, C_{TG}\big|_{\text{parede}}
\label{eq:rw}
\end{equation}

\noindent
em que $C_{TG} = \rho\, Y_{TG}/M_{TG}$ é a concentração molar de triglicerídeo
avaliada junto à superfície, $k_{cat}$ a constante cinética por massa de
catalisador e $w$ a massa de catalisador por área de parede, dada por

\begin{equation}
w = \rho_{wc}\, \delta_{wc}
\label{eq:w}
\end{equation}

A constante cinética foi tomada de \textcite{Melo2023}, que a determinou para
CaO suportado em biochar na metanólise de óleo de soja a \SI{65}{\celsius}:

\begin{equation}
k_{cat} = \SI{5.302e-5}{\meter\cubed\per\kilogram\per\second}
\end{equation}

Na ausência de medidas da camada depositada, adotaram-se
$\rho_{wc} = \SI{1700}{\kilogram\per\meter\cubed}$ e
$\delta_{wc} = \SI{10}{\micro\meter}$, o que resulta em
$w = \SI{1.70e-2}{\kilogram\per\meter\squared}$. Convém registrar que a taxa da
Equação~\eqref{eq:rw} é linear no produto $\rho_{wc}\,\delta_{wc}$, de modo que
a caracterização experimental dessas duas grandezas é a medida de maior retorno
imediato para a acurácia do modelo.

\pendente{substituir $\rho_{wc}$ e $\delta_{wc}$ pelos valores medidos assim que
as caracterizações do \textit{washcoat} estiverem concluídas}

Os fluxos mássicos de cada espécie na parede decorrem da estequiometria da
Equação~\eqref{eq:global}, com sinal negativo para os reagentes e positivo para
os produtos:

\begin{align}
J_{TG}   &= -M_{TG}\, r_w        &&= -872{,}00\, r_w \label{eq:fluxTG}\\
J_{MeOH} &= -3\,M_{MeOH}\, r_w   &&= -\phantom{0}96{,}12\, r_w \label{eq:fluxMeOH}\\
J_{FAME} &= +3\,M_{FAME}\, r_w   &&= +876{,}00\, r_w \label{eq:fluxFAME}\\
J_{GL}   &= +M_{GL}\, r_w        &&= +\phantom{0}92{,}09\, r_w \label{eq:fluxGL}
\end{align}

A soma dos quatro coeficientes vale $-0{,}03$ quilogramas por quilomol de
triglicerídeo consumido, ou \num{3e-5} da massa envolvida, resíduo que decorre
apenas do arredondamento das massas molares. Essa verificação não é
dispensável: coeficientes que não se cancelam injetam ou removem massa total na
parede, e o solver redistribui o desvio pela normalização de $\sum_i Y_i$,
produzindo um erro que não se manifesta em nenhum dos indicadores usuais de
convergência.

\section{Resistência interna ao \textit{washcoat}}

A difusão no interior da camada porosa foi tratada pelo módulo de Thiele para
placa plana, Equações~\eqref{eq:thiele} e~\eqref{eq:efetividade}. Com a
difusividade efetiva no poro tomada da faixa reportada por
\textcite{Fonseca2026} para CaO e CaO/Al$_2$O$_3$, e com
$\delta_{wc} = \SI{10}{\micro\meter}$, obtém-se

\begin{equation}
\phi = 1{,}44
\qquad \Longrightarrow \qquad
\eta = 0{,}62
\end{equation}

Ou seja, pouco menos de \SI{40}{\percent} da atividade nominal do catalisador
depositado é perdida por limitação difusional interna já nessa espessura.

O comportamento da taxa observada com a espessura da camada merece ser
explicitado, por ter consequência direta sobre o procedimento de impregnação.
Substituindo-se a Equação~\eqref{eq:efetividade} na taxa por área de parede,

\begin{equation}
r_{obs} = \eta\, k_v\, \delta_{wc}\, C
        = C \sqrt{k_v D_p}\, \tanh \phi
\label{eq:saturacao}
\end{equation}

\noindent
de modo que a taxa não passa por máximo: ela cresce monotonicamente e satura em
$C\sqrt{k_v D_p}$, valor que independe da espessura depositada. O que se perde
ao engrossar a camada não é taxa, e sim material, que passa a ser depositado sem
contribuir.

Com a calibração adotada, $\phi = \num{1.44}$ a \SI{10}{\micro\meter}
corresponde a $\tanh\phi = \num{0.894}$, ou seja, \SI{89}{\percent} da taxa
máxima alcançável. Aos \SI{15}{\micro\meter} chega-se a \SI{97}{\percent}, e
aos \SI{20}{\micro\meter} a \SI{99}{\percent}. A faixa de trabalho
recomendada é, portanto, de \SIrange{10}{15}{\micro\meter}: abaixo disso
perde-se taxa de forma apreciável, e acima disso acrescenta-se catalisador cujo
retorno é inferior a três pontos percentuais. Essa recomendação foi repassada à
equipe responsável pelas impregnações.

\section{Propriedades físicas e difusividades}
\label{sec:difusividades}

As propriedades adotadas estão reunidas na Tabela~\ref{tab:propriedades}. A
composição de entrada corresponde à razão molar de \num{4}:\num{1} entre metanol
e triglicerídeo,

\begin{equation}
Y_{TG} = \frac{872}{872 + 4 \times 32{,}04} = 0{,}8719,
\qquad Y_{MeOH} = 0{,}1281
\label{eq:composicao}
\end{equation}

\noindent
escolha que merece justificativa, por divergir da razão de \num{6}:\num{1} usual
na prática industrial. A constante cinética adotada na Equação~\eqref{eq:rw} foi
determinada por \textcite{Melo2023} em ensaios conduzidos com
\SI{1}{\percent} em massa de CaO, razão molar metanol:óleo de \num{4}:\num{1} e
\SI{65}{\celsius}. Alimentar o modelo com razão diferente daquela em que a
cinética foi ajustada introduziria inconsistência entre o parâmetro e a condição
de contorno, ainda que o efeito numérico seja pequeno: por se tratar de cinética
de pseudo-primeira ordem, tanto o número de Damköhler quanto a conversão
fracionária são razões e independem, em primeira aproximação, do valor absoluto
da concentração de entrada. Essa independência foi verificada por simulação, e
os valores medidos constam da Seção~\ref{sec:invariancia}.

A escolha tem ainda um benefício colateral. Segundo as medidas de solubilidade
de \textcite{Firth2014}, a \SI{65}{\celsius} o óleo vegetal dissolve metanol até
a razão de \num{3}:\num{1}; uma alimentação de \num{4}:\num{1} está, portanto,
mais próxima do limite de miscibilidade do que uma de \num{6}:\num{1}, o que
atenua, sem eliminar, a fragilidade discutida na
Seção~\ref{sec:limitacoes}.

\begin{table}[htb]
\centering
\caption{Propriedades físicas e condições operacionais adotadas}
\label{tab:propriedades}
\begin{tabular}{l c c l}
\toprule
Propriedade & Símbolo & Valor & Unidade \\
\midrule
Massa específica do meio      & $\rho$      & \num{868}      & \si{\kilogram\per\meter\cubed} \\
Viscosidade dinâmica          & $\mu$       & \num{1.05e-2}  & \si{\pascal\second} \\
Difusividade efetiva          & $D_i$       & \num{7.0e-11}  & \si{\meter\squared\per\second} \\
Temperatura                   & $T$         & \num{65}       & \si{\celsius} \\
Velocidade média              & $\bar{u}$   & \num{1.0e-3}   & \si{\meter\per\second} \\
Altura do canal               & $H$         & \num{1.1}      & \si{\milli\meter} \\
Comprimento do canal          & $L$         & \num{50}       & \si{\milli\meter} \\
Massa específica do washcoat  & $\rho_{wc}$ & \num{1700}     & \si{\kilogram\per\meter\cubed} \\
Espessura do washcoat         & $\delta_{wc}$ & \num{10}     & \si{\micro\meter} \\
\bottomrule
\end{tabular}
\end{table}

A igualdade imposta entre as difusividades das quatro espécies não é uma
simplificação cosmética, e merece ser justificada com algum cuidado. Na
formulação de Fick, o fluxo difusivo de cada espécie é escrito como
$\mathbf{j}_i = -\rho D_i \nabla Y_i$. A conservação da massa total exige que a
soma desses fluxos se anule,

\begin{equation}
\sum_i \mathbf{j}_i = -\rho \sum_i D_i \nabla Y_i = 0
\label{eq:somafluxos}
\end{equation}

\noindent
e, como $\sum_i Y_i = 1$ implica $\sum_i \nabla Y_i = 0$, a
Equação~\eqref{eq:somafluxos} só é satisfeita identicamente quando todas as
difusividades forem iguais entre si. Com difusividades distintas, a soma das
frações mássicas deriva da unidade a cada passo de tempo, e o solver a
restabelece por normalização, distribuindo o desvio proporcionalmente a $Y_i$.
O erro recai, portanto, quase integralmente sobre a espécie majoritária.

No sistema estudado o triglicerídeo responde por \SI{82}{\percent} da massa e o
metanol possui difusividade cerca de noventa vezes maior, dada a diferença de
tamanho molecular. Nessas condições, simulações preliminares conduzidas com
difusividades individuais apresentaram erro de \SI{374}{\percent} no balanço de
massa do triglicerídeo, sem que resíduos, somas de frações mássicas ou fluxos de
parede indicassem qualquer anormalidade. A adoção de um valor efetivo único,
$D_i = \SI{7.0e-11}{\meter\squared\per\second}$, reduziu esse erro a
\SI{0.013}{\percent}.

A perda física associada à simplificação é real e deve ser declarada: a
seletividade difusional entre espécies é suprimida. No presente caso a
simplificação é defensável, pois o reagente limitante é o triglicerídeo, o
metanol se encontra em excesso e o objeto de interesse é a competição entre
cinética de parede e transporte externo do próprio triglicerídeo. Caso a
seletividade entre espécies venha a se tornar relevante, a alternativa correta é
a formulação de Maxwell-Stefan, condicionada à obtenção de difusividades
binárias para todos os pares.

\section{Malha}

Empregou-se malha do tipo \textit{trimmer} com dimensão de base de
\SI{50}{\micro\meter}, igual à espessura do domínio na direção de simetria, o
que força uma única célula nessa direção. Junto às duas paredes foram geradas
camadas prismáticas com os parâmetros da Tabela~\ref{tab:malha}, resultando em
\num{88000} células.

\begin{table}[htb]
\centering
\caption{Parâmetros da malha}
\label{tab:malha}
\begin{tabular}{l c}
\toprule
Parâmetro & Valor \\
\midrule
Dimensão de base                      & \SI{50}{\micro\meter} \\
Número de camadas prismáticas         & \num{15} \\
Espessura total das camadas           & \SI{200}{\micro\meter} \\
Razão de crescimento                  & \num{1.2} \\
Espessura da primeira camada          & \SI{2.8}{\micro\meter} \\
Número total de células               & \num{88000} \\
\bottomrule
\end{tabular}
\end{table}

O refino junto à parede não é opcional neste problema. Com $Sc = \num{1.7e5}$, o
gradiente de fração mássica na superfície catalítica atinge

\begin{equation}
\left. \frac{\partial Y}{\partial n} \right|_{\text{parede}}
= \frac{J}{\rho D} \approx \SI{3.9e3}{\per\meter}
\end{equation}

\noindent
o que corresponde a uma variação de $\Delta Y \approx 0{,}20$ ao longo de uma
célula de \SI{50}{\micro\meter}, ou \SI{40}{\percent} do valor de parede. Sem as
camadas prismáticas, o valor de face lido pela condição de contorno é uma
extrapolação sem significado físico, e o balanço global de massa apresenta
desvios de dezenas de pontos percentuais já nos primeiros segundos de simulação.

\section{Discretização e marcha temporal}

Adotou-se formulação implícita transiente com \num{30} iterações internas por
passo de tempo. O erro de discretização temporal escala linearmente com o passo,
dada a ordem do esquema de Euler implícito; mediu-se contribuição de
aproximadamente \num{4.5} pontos percentuais no balanço de massa com
$\Delta t = \SI{0.05}{\second}$, reduzida a cerca de \num{0.45} ponto com
$\Delta t = \SI{5e-3}{\second}$. A marcha foi conduzida com
$\Delta t = \SI{5e-3}{\second}$ na partida e $\Delta t = \SI{2e-2}{\second}$ no
trecho longo, com refinamento final antes da extração dos resultados.

\section{Procedimento de verificação}
\label{sec:verificacao}

A verificação apoia-se em três testes independentes, todos necessários. A
experiência acumulada ao longo do trabalho mostrou que nenhum deles,
isoladamente, é suficiente para detectar os modos de falha encontrados.

O primeiro é a \textbf{consistência estequiométrica}. A razão entre as taxas de
acúmulo de FAME e de consumo de triglicerídeo no domínio deve igualar a razão
mássica estequiométrica,

\begin{equation}
\frac{\mathrm{d}M_{FAME}/\mathrm{d}t}{\mathrm{d}M_{TG}/\mathrm{d}t}
= \frac{876}{872} = 1{,}004587
\label{eq:esteq}
\end{equation}

\noindent
restrição exata, independente da malha, do passo de tempo e do estado de
convergência. É o único teste que flagra a violação de conservação descrita na
Seção~\ref{sec:difusividades}.

O segundo é a \textbf{independência de malha}, verificada pela reprodução do
resultado central em configurações distintas de discretização.

O terceiro é o \textbf{balanço global de massa} por espécie, escrito como

\begin{equation}
\frac{\mathrm{d}}{\mathrm{d}t} \int_V \rho\, Y_i\, \mathrm{d}V
= \dot{m}\left( Y_{i,\text{ent}} - Y_{i,\text{sai}} \right) + \oint_{A_p} J_i\, \mathrm{d}A
\label{eq:balanco}
\end{equation}

\noindent
com todos os termos medidos de forma independente. É o teste mais trabalhoso e o
mais importante, pois resíduos, somas de frações mássicas e fluxos de parede
podem apresentar-se simultaneamente corretos enquanto o balanço falha por ordens
de grandeza.

Uma observação de ordem prática, decorrente da experiência deste trabalho, deve
acompanhar a Equação~\eqref{eq:balanco}. O termo convectivo não pode ser
avaliado pela média de fração mássica ponderada por vazão na fronteira de saída,
pelas razões documentadas na Seção~\ref{sec:reports}. Em seu lugar, emprega-se a
integração numérica de um perfil transversal medido por sonda de linha, ou
calcula-se a conversão diretamente pela integral do fluxo de parede.

\section{Limitações declaradas do modelo}
\label{sec:limitacoes}

Três hipóteses do modelo merecem ser explicitadas como limitações, por terem
sido confrontadas com dados experimentais da literatura durante a redação deste
texto. São enumeradas a seguir em ordem decrescente de gravidade.

\begin{enumerate}
\itemsep10pt

\item \textbf{Hipótese de fase única.}

O modelo trata o meio reacional como uma única fase líquida homogênea. As
medidas de solubilidade de \textcite{Firth2014} indicam, porém, que a
\SI{65}{\celsius} o óleo vegetal dissolve metanol apenas até a razão molar de
\num{3}:\num{1}, ao passo que a alimentação aqui adotada é de \num{4}:\num{1}.
Na entrada do canal, portanto, parte do álcool pode constituir fase separada, e
o sistema seria bifásico.

O quadro é, na verdade, mais desfavorável do que essa comparação sugere.
\textcite{Muller2019} registra que a mistura de óleo e metanol só se torna
homogênea a partir de \SIrange{15}{20}{\percent} de conversão, patamar muito
acima da conversão por passe de \SI{1.4}{\percent} obtida neste trabalho. O
canal simulado operaria, assim, em regime bifásico em toda a sua extensão.

Dois atenuantes se aplicam, e convém dimensioná-los com honestidade. O primeiro
é que o éster formado atua como cossolvente, de modo que a miscibilidade aumenta
à medida que a reação avança; isso importa pouco em um único passe, mas passa a
importar em configurações com reciclo ou com várias peças em série, nas quais a
conversão acumulada pode ultrapassar aquele patamar. O segundo é que o reagente
limitante e o objeto de interesse deste trabalho é o triglicerídeo, cuja
concentração na fase oleosa é a que governa a taxa de parede, ao passo que o
metanol, estando em excesso, não limita.

Ainda assim, esta é a hipótese mais frágil do modelo, e convém delimitar com
precisão o que ele descreve. O que o modelo resolve com rigor é o transporte de
triglicerídeo no interior da fase oleosa até a parede catalítica. A
disponibilidade de metanol nessa parede, governada pela transferência entre
fases, fica fora de seu alcance. A verificação experimental dessa hipótese
consta das atividades previstas no Capítulo~\ref{cap:cronograma}, e a extensão
do modelo a duas fases é a principal direção de trabalho futuro que dela
decorre.

\item \textbf{Propriedades físicas constantes.}

A massa específica e a viscosidade foram tomadas como constantes.
\textcite{Firth2014} atribui à viscosidade elevada do meio rico em
triglicerídeo uma fase inicial de inibição da reação, o que indicaria
dependência da taxa com a composição.

No presente caso a hipótese é defensável precisamente pelo resultado obtido:
com conversão por passe de \SI{1.4}{\percent}, a composição do meio percorre
faixa estreita ao longo do canal, e as propriedades não chegam a variar de forma
apreciável. A hipótese deixaria de valer em configurações de conversão elevada.

\item \textbf{Cinética irreversível.}

A reação global é reversível, e o modelo a trata como irreversível. A
aproximação é válida enquanto a composição estiver longe do equilíbrio,
condição também assegurada pela baixa conversão por passe.

\end{enumerate}

As três limitações têm a mesma origem e o mesmo remédio. Todas decorrem de se
extrapolar um modelo construído para conversão baixa, e todas serão reexaminadas
quando a calibração experimental prevista no Capítulo~\ref{cap:cronograma}
fornecer dados em faixas mais amplas.
